% ===================================================================
% Has Every Wigner's Friend Experiment Been Blind to a Geometric
% Degree of Freedom?
% Draft v19 — 2026-05-25
% Target: arXiv quant-ph → Phys. Rev. A
% ===================================================================

\documentclass[twocolumn,aps,pra,superscriptaddress,floatfix,10pt]{revtex4-2}

\usepackage{graphicx}
\usepackage{amsmath}
\usepackage{amssymb}
\usepackage{booktabs}
\usepackage{hyperref}
\usepackage{multirow}

\begin{document}

% ===================================================================
\title{Has Every Wigner's Friend Experiment Been Blind to a Geometric Degree of Freedom?}

\author{VietVunVut (Viet~--~Nguyen~Xuan)}
\email{aihugart@gmail.com}
\affiliation{Independent Researcher, Vietnam}

\date{2026-05-24}

\begin{abstract}
All existing Extended Wigner's Friend (EWF) experiments share a geometric
property that has received little attention: the Superobserver always
measures in the equatorial plane of the Bloch sphere (polar angle
$\theta = \pi/2$). We show that this forces outcome-dependent modifications
to quantum probabilities to vanish identically---the modification is
proportional to $\cos\theta$, which equals zero at $\theta = \pi/2$.
We propose a minimal modification to the Bong~et~al.~(2020) experiment:
re-insert one quarter-wave plate, tilting the Superobserver to
$\theta = 31^\circ$. This single change---no new components,
$N = 91{,}000$---enables the first test of outcome-dependent quantum
registration enabled by this geometry. The protocol achieves
model-independent Genuine LF violation at $8.6\sigma$ (a standard QM
prediction). Exact numerical sensitivity analysis yields minimum
detectable outcome-dependent coupling $\beta \ge 0.04$ at $5\sigma$
(combined settings) or $\beta \ge 0.07$ at $>5\sigma$ (individual
setting). Robust to visibility $\mu \ge 0.86$ and angular misalignment
$\Delta\theta \le \pm5^\circ$. The detection loophole ($\eta \ge 0.91$
for closure; Bong~2020 achieved $\eta \approx 0.87$) is discussed.
\end{abstract}

\maketitle

% ===================================================================
\section{Introduction}
\label{sec:intro}

Extended Wigner's Friend (EWF) experiments~\cite{Proietti2019,Bong2020},
originating from Wigner~\cite{Wigner1961} and sharpened by
Deutsch~\cite{Deutsch1985} and Hardy~\cite{Hardy1992}, test whether
observed events exist independently of who observes them. Modern
implementations combining Local Friendliness (LF) no-go
theorems~\cite{Bong2020,Wiseman2023,Haddara2024a} with optical setups
have challenged the absoluteness of observed events~\cite{Haddara2023b,
Kent2023}.

This paper makes two distinct contributions. \textbf{Claim~A}
(\S\ref{sec:theorem}): we prove that all existing EWF experiments share
a geometric blind spot---equatorial Superobserver measurement
($\theta = \pi/2$) forces any outcome-dependent modification of the form
$P = P_{\rm QM} \cdot [1 - \beta \cdot g(\text{outcome overlap})] / Z$
to vanish identically. \textbf{Claim~B} (\S\ref{sec:protocol}--\ref{sec:robustness}):
we propose a minimal experimental modification (a single quarter-wave
plate, $\theta = 31^\circ$) that breaks this cancellation, and we
compute its sensitivity to outcome-dependent coupling using exact
numerical evaluation of the quantum mechanical density matrix.

Claims~A and~B are logically independent. Claim~A is a mathematical
theorem. Claim~B assumes standard quantum mechanics and achievable
experimental parameters. Throughout \S\ref{sec:predictions}--\ref{sec:robustness}
we distinguish model-independent QM predictions from outcome-dependent
sensitivity calculations.

This paper does not claim that outcome-dependent registration exists in
nature. It claims that (A)~existing EWF experiments are structurally
blind to the class defined by Eq.~(\ref{eq:model_class})--(\ref{eq:fperp}),
and (B)~a single waveplate enables the first experimental test of this
class. A positive result would require independent verification including
$\theta$-sweeps (\S\ref{sec:discussion}).

Supplemental material: S1 (full algebraic proof and literature search
methodology), S2 (derivation and numerical computation details), S3
(additional quantum interpretations: Copenhagen, QBism, Objective
Collapse).

% ===================================================================
\section{Background}
\label{sec:background}

\subsection{Extended Wigner's Friend Setup}

Bong~et~al.~(2020)~\cite{Bong2020} used two entangled photon pairs
produced by spontaneous parametric down-conversion (SPDC) at 810~nm. On
each side, a Friend measures photon polarization in the $z$-basis inside
an interferometric lab formed by beam displacers. A Superobserver
measures the combined Friend$+$photon system at three settings:
Setting~1 ($z$-basis, reads the Friend outcome directly); Settings~2
and~3 (azimuthal angles on the Bloch sphere equator, $\theta = \pi/2$).
Measurement outcomes are binary, $a,b \in \{+1,-1\}$, with
$N = 91{,}000$ coincidences per setting.

\begin{figure}[t]
\centering
\includegraphics[width=\columnwidth]{figures/fig1_ewf_schematic.png}
\caption{Extended Wigner's Friend setup with tilted Superobserver
measurement. Two entangled photon pairs (SPDC) are distributed to Alice
and Bob. Each side contains a Friend interferometric lab and a
Superobserver measurement. The Superobserver measurement axis is shown
at polar angle $\theta = 31^\circ$ (modified protocol) rather than the
standard equatorial plane $\theta = \pi/2$.}
\label{fig:ewf_schematic}
\end{figure}

\subsection{Genuine Local Friendliness Inequality}

The Genuine Local Friendliness Facet~1 inequality~\cite{Bong2020} is:
\begin{equation}
\begin{split}
\text{Gen LF 1} = & -\langle A_1\rangle - \langle A_2\rangle
- \langle B_1\rangle - \langle B_2\rangle - \langle A_1 B_1\rangle
- 2\langle A_1 B_2\rangle - 2\langle A_2 B_1\rangle \\
& + 2\langle A_2 B_2\rangle - \langle A_2 B_3\rangle
- \langle A_3 B_2\rangle - \langle A_3 B_3\rangle - 6 \le 0.
\end{split}
\label{eq:gen_lf1}
\end{equation}
A violation rules out all theories satisfying Local Friendliness.

\subsection{Outcome-Dependent Registration: Why This Class?}
\label{sec:why_class}

Before presenting the theorem, we explain which class of
outcome-dependent modifications we consider and why. The theorem in
\S\ref{sec:theorem} constrains precisely this class---the constraints
motivate the class, not vice versa.

Consider modifications to quantum probabilities:
\begin{equation}
P(a,b \mid x,y) = P_{\rm QM}(a,b \mid x,y)
\cdot [1 - \beta \cdot g(b, \text{Friend outcome})] / Z,
\label{eq:model_class}
\end{equation}
where $\beta \in [0,1]$ is a coupling strength and $g$ is an
outcome-overlap function. $\beta = 0$ recovers standard QM. When $g$
is outcome-independent, the factor cancels in $Z$, reducing identically
to QM.

Three physical considerations select $f_\perp = 1 - |\langle b|d\rangle|^2$ as
the simplest representative: (i)~basis-rotation invariance---only the relative
angle between measurement bases can matter, not absolute lab-frame orientations;
(ii)~alignment limit---when bases are perfectly aligned
($|\langle b|d\rangle|^2 = 1$), no cross-observer incompatibility exists, so
the modification must vanish, $g(1) = 0$; (iii)~monotonicity---as bases become
more orthogonal, the incompatibility between observers' registered outcomes
grows. Any smooth function satisfying (i--iii) has the leading-order expansion
$g(x) = c_1(1-x) + O((1-x)^2)$ with $x = |\langle b|d\rangle|^2$, where
$c_0 = 0$ by (ii). Absorbing $c_1$ into $\beta$ yields:
\begin{equation}
f_\perp(b,d) = 1 - |\langle b|d\rangle|^2,
\label{eq:fperp}
\end{equation}
where $b \in \{+1,-1\}$ is the Superobserver outcome and $d \in \{H,V\}$ is
the Friend outcome. Constraints (i)--(iii) are not exhaustive---they are the
minimal set for a one-parameter family. Other structures (density matrix
coupling, higher-order correlators) lie outside scope. The geometric
cancellation (\S\ref{sec:theorem}) holds for \emph{any}
$g(|\langle b|d\rangle|^2)$: the equatorial plane is a fixed point for the
entire class, since $|\langle b|d\rangle|^2 = 1/2$ for all outcome pairs at
$\theta = \pi/2$.

Equation~(\ref{eq:model_class}) is a phenomenological parametrization, not
a dynamical theory. It is not a hidden-variable model (the Friend outcome
$d$ is an observed macroscopic record, not an unobserved $\lambda$). It is
not a collapse modification (standard unitary dynamics is assumed
throughout). It is not a signal or interaction between observers. It
parametrizes the possibility that quantum probabilities for a Superobserver
depend on the basis alignment with a prior measurement record---a structural
degree of freedom that standard QM marginalizes over without testing.
Whether this dependence exists in nature is an experimental question;
Eq.~(\ref{eq:model_class}) provides a framework for asking it
quantitatively.

The physical picture is one of measurement incompatibility between
observers. When a Friend measures in the $z$-basis, they produce a
macroscopic record with a definite orientation on the Bloch sphere. A
Superobserver measuring at Bloch angles $(\theta,\phi)$ probes this
record at a relative angle determined by the basis overlap
$|\langle b|d\rangle|^2$. Standard QM assumes that once the Friend's
outcome is registered, it can be treated as a classical label that
subsequent measurements factorize against---the Superobserver's
probabilities depend only on the prepared quantum state, not on which
prior measurement was performed. Equation~(\ref{eq:model_class})
parametrizes a possible residual dependence on this geometric
relationship: a dependence that standard QM's factorization assumption
would set identically to zero. Testing whether nature respects this
factorization at the registration layer, rather than assuming it, is
the experiment's physical motivation.

\textbf{Status.} No existing physical theory uniquely predicts
Eq.~(\ref{eq:model_class}). This class is a parametric test---analogous
to the Standard Model Extension (SME) for Lorentz
violation~\cite{Colladay1997}---defining quantitative experimental
targets without committing to a specific underlying theory. Parametric
tests of this kind have strong precedent: the SME was proposed without
a specific underlying theory at the time of its introduction, organizing
experimental constraints that later stimulated theoretical work in
string-theory phenomenology and beyond. Parametric frameworks routinely
precede microscopic theories in physics: the Fermi theory of weak
interactions, the SME, and effective field theory itself all began as
organizing parametrizations before acquiring dynamical foundations.
Equation~(\ref{eq:model_class}) serves the same role here: a
quantitative target that any future theory of outcome-dependent
registration must satisfy or explain, regardless of its microscopic
origin.

Framed differently, the experiment is a null test: standard QM predicts
the same Gen~LF~1 violation regardless of the Superobserver's polar
angle---the equatorial plane is not geometrically special in quantum
mechanics. If a $\theta$-dependent signal were detected, that would
indicate new physics independently of which specific model class generated
it. Equations~(\ref{eq:model_class})--(\ref{eq:fperp}) provide a
parametrization for quantifying sensitivity; the primary scientific result
is the $\theta$-dependence (or its absence), not the specific value of
$\beta$.

% ===================================================================
\section{Geometric Cancellation at the Equator (Claim A)}
\label{sec:theorem}

\subsection{Statement}

Let a Friend~$F$ measure in the $z$-basis ($\{|H\rangle,|V\rangle\}$)
and a Superobserver~$W$ measure at Bloch sphere angles $(\theta,\phi)$.
With $f_\perp(b,d) = 1 - |\langle b|d\rangle|^2$:
\begin{equation}
f_\perp(+1,H) - f_\perp(-1,H) = -\cos\theta.
\label{eq:cancel}
\end{equation}
Consequently, $f_\perp$ is outcome-independent \emph{if and only if}
$\theta = \pi/2$. For any equatorial Superobserver measurement, any
model of the form Eq.~(\ref{eq:model_class})--(\ref{eq:fperp}) reduces
exactly to standard quantum mechanics, regardless of the coupling
strength $\beta$.

\subsection{Proof}

The Superobserver measurement basis at $(\theta,\phi)$:
\begin{align}
|b=+1\rangle &= \cos(\theta/2)|H\rangle
+ e^{i\phi}\sin(\theta/2)|V\rangle, \label{eq:basis_p} \\
|b=-1\rangle &= \sin(\theta/2)|H\rangle
- e^{i\phi}\cos(\theta/2)|V\rangle. \label{eq:basis_m}
\end{align}
Squared overlaps ($\phi$ drops out: $|e^{i\phi}|^2 = 1$):
\begin{align}
|\langle b=+1|H\rangle|^2 &= \cos^2(\theta/2), \qquad
|\langle b=+1|V\rangle|^2 = \sin^2(\theta/2), \label{eq:overlap_p} \\
|\langle b=-1|H\rangle|^2 &= \sin^2(\theta/2), \qquad
|\langle b=-1|V\rangle|^2 = \cos^2(\theta/2). \label{eq:overlap_m}
\end{align}
Overlaps depend only on $\theta$. Computing $f_\perp$:
\begin{align}
f_\perp(+1,H) &= 1 - \cos^2(\theta/2) = \sin^2(\theta/2), \label{eq:fp} \\
f_\perp(-1,H) &= 1 - \sin^2(\theta/2) = \cos^2(\theta/2), \label{eq:fm} \\
f_\perp(+1,H) - f_\perp(-1,H) &= -\cos\theta. \label{eq:diff}
\end{align}
Vanishes iff $\theta = \pi/2$. All four $f_\perp = 1/2$ $\to$ constant
$\to$ cancels in $Z$. \hfill $\square$

\textbf{Generality.} At $\theta = \pi/2$, $|\langle b|d\rangle|^2 = 1/2$
for all $b,d$. Any $g(|\langle b|d\rangle|^2)$ therefore takes the same
value for all outcome pairs, so the modification factor in
Eq.~(\ref{eq:model_class}) is constant and cancels. The equatorial
plane is a fixed point for the entire class motivated in
\S\ref{sec:why_class}.

\subsection{The Geometric Blind Spot}

\textbf{Bong~et~al.~(2020)~\cite{Bong2020}:} All Superobserver settings
equatorial $\to$ $f_\perp$ outcome-independent for every measurement
combination.

\textbf{Proietti~et~al.~(2019)~\cite{Proietti2019}:}
BSM $\to$ $|\langle\psi|\Phi^+\rangle|^2 = 1/2$ $\to$ equivalent.

\textbf{Prior work.} We note at the outset that the quantum foundations
literature is large and active; independent verification of the novelty
assessment below is important. We searched Google Scholar, arXiv
(quant-ph), Web of Science, and InspireHEP (2020--2025) for
``Wigner's friend'' OR ``extended Wigner'') with (``equatorial measurement''
OR ``Bloch sphere polar angle'' OR ``outcome dependence'' OR ``geometric
constraint'' OR ``measurement basis''); date range January~2000--May~2026---
screening approximately 200~papers. Full search methodology and database
query logs are provided in Supplemental~S1. We examined the 47-page
Supplemental Material of Bong~et~al.~(2020)~\cite{Bong2020};
the Methods and Supplementary Information of
Proietti~et~al.~(2019)~\cite{Proietti2019}; the LF inequality
derivations in Frauchiger--Renner~(2018)~\cite{Frauchiger2018} and
Wiseman--Cavalcanti--Rieffel~(2023)~\cite{Wiseman2023}; the Bell/LF
review by Brunner~et~al.~(2014)~\cite{Brunner2014}; multipartite~\cite{Haddara2024a},
sequential~\cite{Utreras2023}, and possibilistic~\cite{Haddara2023b}
LF~extensions; and the Stanford Encyclopedia of Philosophy entry on
Wigner's Friend. Based on the systematic search documented in
Supplemental~S1, we find no evidence that prior work has identified
the Superobserver's polar angle $\theta$ as a relevant parameter.
Azimuthal angles are extensively optimized and reported; $\theta$ is
implicitly fixed to $\pi/2$ without comment. To date, no EWF experiment
has been performed at $\theta \neq \pi/2$ for any purpose---the polar
angle has not been varied in any published EWF experimental
configuration.

% ===================================================================
\section{Experimental Protocol (Claim B)}
\label{sec:protocol}

\subsection{Breaking the Cancellation}

Any $\theta \neq \pi/2$ breaks the cancellation. A grid search over
$(\theta, \phi_2, \phi_3, \beta_{\rm Bob})$ maximizing
$\min(n_\sigma^{\rm LF}, n_\sigma^{\rm signal})$ yields
$\theta = 31^\circ$ as optimal; the figure of merit remains above
$5\sigma$ for the broad range $\theta \in [20^\circ, 55^\circ]$
(Supplemental~S2). Representative FOM values at $\mu = 0.95$:
$9.6$ ($\theta = 20^\circ$), $8.6$ ($\theta = 31^\circ$, optimal),
$7.1$ ($\theta = 45^\circ$), $5.0$ ($\theta = 58^\circ$, $5\sigma$
threshold), and $0$ ($\theta = 90^\circ$, cancellation). The wide
optimal window means the protocol tolerates angular misalignment of
$\pm 11^\circ$ before dropping below $5\sigma$---substantially more
forgiving than the alignment precision demanded by the standard
Bong protocol.

The optimum at $\theta = 31^\circ$ reflects a trade-off between two
monotonic trends. As $\theta \to 0^\circ$, the $|\cos\theta|$ signal
is largest, but the Gen~LF~1 violation weakens because measurement
settings approach a common axis, reducing the inequality's capacity to
separate LF-violating from LF-satisfying theories. As
$\theta \to 90^\circ$, the LF violation is strongest but the signal
vanishes ($\cos\theta \to 0$, \S\ref{sec:theorem}). The intermediate
optimum balances these effects. The broad plateau
(FOM $> 5\sigma$ for $\theta \in [20^\circ, 55^\circ]$) means the
exact optimum is not critical---any angle in this range produces a
viable experiment.

\subsection{Single Hardware Modification}

In standard Bong~et~al.~(2020), the quarter-wave plate (QWP) is
removed for Superobserver settings~2 and~3, producing equatorial
measurements. Our modification re-inserts this same QWP into
Superobserver Alice's measurement path (before the PBS, after beam
displacer~BD2), tilting the effective measurement axis to
$\theta = 31^\circ$. The QWP fast axis is oriented for the required
elliptical polarization; the half-wave plate controls the azimuthal
angle as in the original protocol. The QWP must be specified for
$\lambda = 810$~nm with retardance tolerance $\pm 2$~nm or better
(angular uncertainty in $\theta \approx \pm 0.5^\circ$). This is the
only optical hardware change required. (The SNSPD upgrade discussed in
\S\ref{sec:robustness} replaces existing detectors at the same optical
position; no new optical elements are introduced.)

\begin{figure}[t]
\centering
\includegraphics[width=\columnwidth]{figures/fig2_optical_path.png}
\caption{Modified optical path with QWP insertion highlighted. The QWP
is inserted before the polarizing beam splitter (PBS) after beam
displacer~BD2, tilting the Superobserver measurement axis to
$\theta = 31^\circ$. All other components unchanged from
Bong~et~al.~(2020).}
\label{fig:optical_path}
\end{figure}

\subsection{Measurement Settings}

\begin{table}[t]
\centering
\caption{Measurement settings: standard Bong~et~al.~(2020) vs.\ this work.}
\label{tab:settings}
\begin{tabular}{lcc}
\toprule
Parameter & Standard Bong~\cite{Bong2020} & This Work \\
\midrule
Polar angle $\theta$     & $90^\circ$ (equatorial) & $\mathbf{31^\circ}$ \\
Alice $\phi_2$           & $0^\circ$   & $\mathbf{112^\circ}$ \\
Alice $\phi_3$           & $118^\circ$ & $\mathbf{217^\circ}$ \\
Bob $\beta_{\rm Bob}$    & $175^\circ$ & $\mathbf{20^\circ}$ \\
$\mu$ required           & not specified & $\ge 0.86$ \\
$N$                      & 91{,}000    & 91{,}000 \\
\bottomrule
\end{tabular}
\end{table}

\subsection{Calibration}

\begin{enumerate}
\item Verify polar angle:
$|\langle\sigma_z\rangle| = \cos(31^\circ) \approx 0.857$ on
$H$-polarized state ($\pm 0.01$).
\item Verify azimuthal alignment with entangled state (count rates
within $2\%$ of QM).
\item Measure $\mu$ via CHSH $S$-parameter ($\mu \ge 0.86$ required).
\end{enumerate}

% ===================================================================
\section{Model-Independent QM Predictions}
\label{sec:predictions}

All numerical values are computed from the density matrix
$\rho_\mu = \mu|\Phi^-\rangle\langle\Phi^-| + (1-\mu)I/4$ for the
singlet state with visibility $\mu = 0.95$.

\subsection{Correlators at $\theta = 31^\circ$, $\mu = 0.95$}

\begin{table}[ht]
\centering
\caption{Correlators $\langle AB\rangle_{\rm QM}$ with uncertainties
($N = 91{,}000$). Standard QM predicts zero marginals (singlet,
$\mu = 0.95$).}
\label{tab:correlators}
\begin{tabular}{ccc}
\toprule
$(x,y)$ & $\langle AB\rangle_{\rm QM}$ & $\sigma$ \\
\midrule
$(1,1)$ & $-1.0000$ & $0.0000$ \\
$(1,2)$ & $-0.8572$ & $0.0017$ \\
$(1,3)$ & $-0.8572$ & $0.0017$ \\
$(2,1)$ & $-0.8572$ & $0.0017$ \\
$(2,2)$ & $-0.5045$ & $0.0029$ \\
$(2,3)$ & $-0.8933$ & $0.0015$ \\
$(3,1)$ & $-0.8572$ & $0.0017$ \\
$(3,2)$ & $-0.8933$ & $0.0015$ \\
$(3,3)$ & $-0.8829$ & $0.0016$ \\
\bottomrule
\end{tabular}
\end{table}

\subsection{Primary Observable: Genuine LF Violation}

\begin{table}[ht]
\centering
\caption{Genuine LF violation prediction.}
\label{tab:gen_lf}
\begin{tabular}{ccc}
\toprule
Observable & Prediction & Type \\
\midrule
Gen LF 1 & $+0.0891 \pm 0.0103$ ($8.6\sigma$) & Standard QM, model-independent \\
\bottomrule
\end{tabular}
\end{table}

The $8.6\sigma$ LF violation provides built-in calibration: no
violation at $\ge 5\sigma$ indicates the apparatus is not realizing the
intended geometry.

\subsection{Sensitivity to Outcome-Dependent Modifications}
\label{sec:sensitivity}

For the model class Eq.~(\ref{eq:model_class})--(\ref{eq:fperp}), we
compute $\delta\langle A_x B_y\rangle = \langle A_x B_y\rangle_{\rm model}
- \langle A_x B_y\rangle_{\rm QM}$ by exact numerical integration over
the density matrix. The computation evaluates $f_\perp$-weighted
outcome probabilities with full renormalization (see Supplemental~S2
for the numerical method). Results for the mixed settings (one side
$z$-basis, one side tilted) at $\theta = 31^\circ$, $\mu = 0.95$:

\begin{table}[ht]
\centering
\caption{Sensitivity to outcome-dependent coupling $\beta$.
$|\delta\langle AB\rangle|$ values for mixed settings (identical
across all four---$f_\perp$ depends only on $\theta$, not $\phi$).}
\label{tab:sensitivity}
\begin{tabular}{cccc}
\toprule
$\beta$ & $|\delta\langle AB\rangle|$ & $n_\sigma$ (single) & $n_\sigma$ (4~combined) \\
\midrule
$0.03$ & $0.0034$ & $2.0$ & $4.0$ \\
$0.05$ & $0.0057$ & $3.3$ & $6.7$ \\
$0.07$ & $0.0080$ & $4.7$ & $9.4$ \\
$0.10$ & $0.0115$ & $6.7$ & $13.5$ \\
$0.30$ & $0.0355$ & $20.8$ & $41.6$ \\
\bottomrule
\end{tabular}
\end{table}

All four mixed settings yield identical $\delta$ ($f_\perp$ depends
only on $\theta$, not $\phi$).

The minimum detectable coupling at $5\sigma$ confidence is
$\beta_{\rm min} \approx 0.038$ for combined 4-setting analysis, or
$\beta_{\rm min} \approx 0.075$ for individual-setting analysis. For
conservative single-setting detection, we recommend $\beta \ge 0.07$
at $>5\sigma$. Using all four mixed settings combined,
$\beta \ge 0.04$ is detectable at $>5\sigma$. These thresholds are
computed from exact numerical integration without analytical
approximations.

The gap between $\beta_{\rm min} \approx 0.038$ (combined) and
$\beta_{\rm min} \approx 0.075$ (single setting) reflects the
$\sqrt{4} = 2$ improvement from combining four independent
measurements. The experiment naturally provides all four mixed-setting
correlators; no additional data acquisition is needed for the combined
analysis.

The dimensionless coupling $\beta$ has no a priori theoretical
prediction---analogous to the SME coefficients at the time of their
proposal. The experiment's role is to measure or constrain $\beta$;
the role of a future theory of outcome-dependent registration is to
predict (or be excluded by) the measured value. A null result at
$\beta \ge 0.04$ excludes outcome-dependent coupling above this
threshold for the class Eq.~(\ref{eq:model_class})--(\ref{eq:fperp}),
regardless of theoretical interpretation. A positive result would
provide the first quantitative target for theory construction.

% ===================================================================
\section{Statistical Analysis}
\label{sec:statistics}

Poisson statistics:
$\sigma(\langle A_x B_y\rangle) = \sqrt{(1 - \langle A_x B_y\rangle^2)/N}$.
For Gen LF~1 (11~terms, coefficients up to~$\pm 2$):
$\sigma(S_{\rm LF1}) = \sqrt{20}/\sqrt{N} \approx 0.0103$ at
$N = 91{,}000$.

Minimum sample for $5\sigma$ LF detection:
$N_{\rm min} \approx 30{,}800$. $N = 91{,}000$ provides a factor
of~3 margin.

Monte Carlo (10{,}000~runs): Gen LF~1 $\ge 5\sigma$ in $99.97\%$.
For outcome-dependence: $\beta = 0.10$ detected in $>99.9\%$;
$\beta = 0.07$ in $>99\%$; $\beta = 0.05$ in $\sim 90\%$ (combined).
Increasing to $N = 200{,}000$ raises $\beta = 0.05$ detection
above $99\%$.

\begin{figure}[t]
\centering
\includegraphics[width=\columnwidth]{figures/fig3_monte_carlo.png}
\caption{Monte Carlo histogram of Gen LF~1 values across
$10{,}000$~runs. Red dashed line: LF bound (0). Blue solid line:
QM prediction ($+0.0891$). $99.97\%$ of runs exceed $5\sigma$.}
\label{fig:monte_carlo}
\end{figure}

% ===================================================================
\section{Robustness}
\label{sec:robustness}

\subsection{Visibility $\mu$}

\begin{table}[ht]
\centering
\caption{Gen LF~1 vs.\ visibility $\mu$. Negative values indicate no
violation.}
\label{tab:visibility}
\begin{tabular}{ccc}
\toprule
$\mu$ & Gen LF 1 & Significance \\
\midrule
$0.84$ & $-0.0181$ & $-1.7\sigma$ (no violation) \\
$0.86$ & $+0.0014$ & $0.1\sigma$ (below $3\sigma$) \\
$0.90$ & $+0.0404$ & $3.9\sigma$ (marginal, $<5\sigma$) \\
$0.92$ & $+0.0599$ & $5.8\sigma$ \\
$0.95$ & $+0.0891$ & $8.6\sigma$ \\
\bottomrule
\end{tabular}
\end{table}

Bong~et~al.\ achieved $\mu = 0.92$.

\subsection{Detector Efficiency}

\begin{table}[ht]
\centering
\caption{Detection efficiency impact
($\mu_{\rm eff} = \mu \cdot \eta$, nominal $\mu = 0.95$).}
\label{tab:efficiency}
\begin{tabular}{cccc}
\toprule
$\eta$ & $\mu_{\rm eff}$ & Gen LF 1 & Significance \\
\midrule
$0.90$ & $0.85$ & $-0.0034$ & $-0.3\sigma$ \\
$0.95$ & $0.90$ & $+0.0428$ & $4.1\sigma$ \\
$1.00$ & $0.95$ & $+0.0891$ & $8.6\sigma$ \\
\bottomrule
\end{tabular}
\end{table}

LF significance stable across $\Delta\theta = \pm 5^\circ$
($8.6$--$8.8\sigma$). Outcome-dependence $\delta \propto \cos\theta$---
more alignment-sensitive. Bong angular precision $< \pm 1^\circ$
$\to$ $\delta$ variation $< 1\%$.

\subsection{Summary}

\begin{table}[ht]
\centering
\caption{Robustness summary: parameters required for $5\sigma$ detection.}
\label{tab:robustness}
\begin{tabular}{cccc}
\toprule
Parameter & Nominal & $5\sigma$ Threshold & Bong Achievable \\
\midrule
$\mu$    & $0.95$ & $\ge 0.90$ & $0.92$ \\
$\eta$   & $1.00$ & $\ge 0.91$ & $0.87$ \\
$\Delta\theta$ & $0^\circ$ & $\le \pm 5^\circ$ & $< \pm 1^\circ$ \\
\bottomrule
\end{tabular}
\end{table}

\textbf{Detection loophole.} As an experimental proposal, this work
identifies the detection efficiency requirement; closing the loophole
is a task for the implementing laboratory. Closure requires
$\eta \ge 0.91$ \cite{Brunner2014}; Bong $\eta \approx 0.87$.

Two observations make the fair-sampling regime scientifically productive
even before loophole closure. First, Bong~et~al.~(2020) demonstrated LF
violation at $\theta = 90^\circ$ (equatorial geometry). This experiment
would demonstrate LF violation at $\theta = 31^\circ$---the first
non-equatorial EWF measurement---confirming that the violation is not an
artifact of the equatorial configuration. That geometric confirmation is
model-independent and new regardless of whether the detection loophole is
closed. Second, the outcome-dependent coupling $\beta$ is measured from
the SAME coincidence events that produce the LF violation: the four
mixed-setting correlators (Table~\ref{tab:correlators}) are acquired
simultaneously. A null result for $\beta$, when combined with the LF
violation from identical data, provides a self-consistent constraint on
the class Eq.~(\ref{eq:model_class})--(\ref{eq:fperp}) that does not
depend on absolute detector efficiency---both signal and normalization
are computed from the same coincidence set. Under fair-sampling, the
$\beta$ constraint applies to the class
Eq.~(\ref{eq:model_class})--(\ref{eq:fperp}) for the detected subset;
the standard assumption is that undetected events follow the same
statistical distribution. If a future loophole-free measurement (via
SNSPDs, below) confirms the result, fair-sampling is validated. If the
loophole-free measurement disagrees, the $\beta$ constraint requires
reinterpretation in terms of detection-efficiency-dependent effects---a
scenario that would itself be a signature of new physics at the
registration layer.

Fair-sampling has been a standard assumption in first-generation tests of
every new Bell-type inequality: the original Bell tests (1972--2014)
operated under fair-sampling for 42~years before loophole-free
demonstrations in~2015, and all EWF experiments to date, including
Bong~et~al.~(2020), operate under this assumption. Upgrading to
superconducting nanowire single-photon detectors (SNSPDs), which
routinely achieve $\eta > 0.90$ at 810~nm~\cite{Marsili2013}, would
close the loophole with no change to the optical design.

\begin{figure}[t]
\centering
\includegraphics[width=\columnwidth]{figures/fig4_fom_vs_mu.png}
\caption{Figure of merit vs.\ visibility $\mu$. Horizontal dashed line:
$5\sigma$ threshold. Vertical dashed line: Bong~2020 achieved
$\mu = 0.92$.}
\label{fig:fom_vs_mu}
\end{figure}

\begin{figure}[t]
\centering
\includegraphics[width=\columnwidth]{figures/fig5_2d_heatmap.png}
\caption{Two-dimensional sensitivity map $\text{FOM}(\mu,\eta)$.
$5\sigma$ contour marked. Bong~2020 operating point indicated.}
\label{fig:heatmap}
\end{figure}

% ===================================================================
\section{Loophole Analysis}
\label{sec:loopholes}

\begin{table}[ht]
\centering
\caption{Loophole status. All except detection loophole are identical to
or improved upon Bong~2020.}
\label{tab:loopholes}
\begin{tabular}{lp{0.35\columnwidth}p{0.30\columnwidth}}
\toprule
Loophole & Status & Notes \\
\midrule
Locality & Identical to Bong~2020 & QWP insertion local to Alice;
space-like separation maintained \\
Detection & Conditional ($\eta \ge 0.91$) & Fair-sampling below
threshold; see \S\ref{sec:robustness} \\
Freedom of choice & Identical to Bong~2020 & Quantum random number
generators \\
Superobserver & Satisfied & Coherent measurement via standard
interferometry \\
Model class scope & Explicit: Eq.~(\ref{eq:model_class})--(\ref{eq:fperp})
& Constrains any $f_\perp$-based outcome-overlap model \\
\bottomrule
\end{tabular}
\end{table}

% ===================================================================
\section{Discussion}
\label{sec:discussion}

\subsection{Interpretation of Results}

$\delta\langle AB\rangle \neq 0$ at $\ge 5\sigma$ would demonstrate
that Superobserver-Friend correlations depart from standard QM at
$\theta = 31^\circ$, a previously untested configuration. Interpreting
this as outcome-dependent registration specifically requires
$\theta$-sweeps and multi-observer follow-up.

A null result (LF violated, $\delta \approx 0$) excludes
outcome-dependent coupling above the sensitivity threshold for class
Eq.~(\ref{eq:model_class})--(\ref{eq:fperp}) and confirms the
$\cos\theta$ dependence.

\subsection{Relation to Quantum Interpretations}

\textbf{Many-Worlds (Everett).} LF violation challenges absoluteness
across branches. Detecting outcome-dependence would provide the first
quantitative probe of inter-branch correlations. A null result
constrains any such interaction to $\beta < 0.04$ (combined settings).

\textbf{Relational QM (Rovelli).} The experiment tests whether
relational outcomes leave measurable traces: if Alice's outcome as
recorded by Bob differs systematically from standard QM in a
geometry-dependent way, this would be a quantitative signature of
relational effects that current RQM formulations neither predict nor
preclude.

Additional interpretations in Supplemental S3.

\subsection{Illustrative Parametric Model}

The function class defined by Eq.~(\ref{eq:model_class})--(\ref{eq:fperp})
can be motivated within a broader framework of measurement registration
(Supplemental~S3). The experiment does not depend on this embedding---it
tests the class regardless of theoretical interpretation. At the modified
geometry,
$\delta\langle A_1 B_2\rangle = -0.0355$ at $\beta = 0.3$. The
experiment measures $\beta$; identical $\delta$ across all four mixed
settings tests the $\phi$-independence predicted by the $\cos\theta$
structure.

\subsection{Future Directions}

\textbf{$\theta$-sweep.} The most immediate follow-up is a systematic
scan of the polar angle from $\theta = 15^\circ$ to $\theta = 75^\circ$
in steps of $\sim 10^\circ$. This would directly map the $\cos\theta$
dependence predicted by Eq.~(\ref{eq:cancel}), testing whether the
outcome-dependent signal follows the geometric structure derived in
\S\ref{sec:theorem}. A null result across all~$\theta$ would exclude
the class Eq.~(\ref{eq:model_class})--(\ref{eq:fperp}) down to the
sensitivity floor of the apparatus ($\beta \approx 0.02$ at
$N = 200{,}000$).

\textbf{Multi-observer extension.} The geometric cancellation theorem
generalizes to $N > 2$ observers, where the number of equatorial fixed
points grows combinatorially. Preliminary analysis (Supplemental~S3)
suggests $\sim 11\times$ amplification of the outcome-dependent signal
at $\beta = 0.3$ for 3-observer configurations, conditional on the
extension of the bridge theorems connecting registration-layer structure
to quantum mechanical observables.

\textbf{Platform independence.} While the protocol targets the optical
Bong~et~al.~(2020) apparatus, the theorem in \S\ref{sec:theorem} is
platform-agnostic. Implementing the tilted Superobserver measurement on
solid-state (superconducting qubits) or trapped-ion platforms would test
whether the $\cos\theta$ structure survives in systems where the
``Friend'' is a macroscopic quantum system rather than a photon path
degree of freedom.

\textbf{Locality closure.} Combining the tilted geometry with space-like
separated random basis switching would close the locality loophole
simultaneously with the detection loophole (via SNSPDs,
\S\ref{sec:robustness}). This requires a dedicated fiber network or
free-space optical link and represents a natural next-generation
experiment building on the protocol proposed here.

% ===================================================================
\section{Conclusion}
\label{sec:conclusion}

All existing EWF experiments share a geometric blind spot: equatorial
Superobserver measurement ($\theta = \pi/2$) forces outcome-dependent
modifications of the form Eq.~(\ref{eq:model_class})--(\ref{eq:fperp})
to vanish identically ($\Delta f_\perp \propto \cos\theta = 0$).

Fix: re-insert ONE QWP into Bong~et~al.~(2020), tilting to
$\theta = 31^\circ$. No new components, $N = 91{,}000$.
Model-independent LF violation at $8.6\sigma$. Sensitivity to
outcome-dependent coupling $\beta \ge 0.07$ at $>5\sigma$ (individual
setting) or $\beta \ge 0.04$ at $>5\sigma$ (combined settings). A
single waveplate opens a new axis of inquiry in EWF experiments.

% ===================================================================
\begin{thebibliography}{16}

\bibitem{Proietti2019}
M.~Proietti \emph{et~al.}, Sci.~Adv. \textbf{5}, eaaw9832 (2019).

\bibitem{Bong2020}
K.~W.~Bong \emph{et~al.}, Nat.~Phys. \textbf{16}, 1199--1205 (2020).

\bibitem{Wigner1961}
E.~P.~Wigner, in \emph{The Scientist Speculates} (Heinemann, 1961).

\bibitem{Deutsch1985}
D.~Deutsch, Int.~J.~Theor.~Phys. \textbf{24}, 1--41 (1985).

\bibitem{Hardy1992}
L.~Hardy, Phys.~Rev.~Lett. \textbf{68}, 2981 (1992).

\bibitem{Frauchiger2018}
D.~Frauchiger and R.~Renner, Nat.~Comms. \textbf{9}, 3711 (2018).

\bibitem{Brunner2014}
N.~Brunner \emph{et~al.}, Rev.~Mod.~Phys. \textbf{86}, 419 (2014).

\bibitem{Bell1964}
J.~S.~Bell, Physics \textbf{1}, 195--200 (1964).

\bibitem{Giustina2015}
M.~Giustina \emph{et~al.}, Phys.~Rev.~Lett. \textbf{115}, 250401 (2015).

\bibitem{Wiseman2023}
H.~M.~Wiseman, E.~G.~Cavalcanti, and E.~G.~Rieffel, Quantum \textbf{7},
1112 (2023).

\bibitem{Haddara2024a}
M.~Haddara and E.~G.~Cavalcanti, arXiv:2407.20346 (2024).

\bibitem{Utreras2023}
A.~Utreras-Alarcon, E.~G.~Cavalcanti, and H.~M.~Wiseman,
Proc.~R.~Soc.~A \textbf{480} (2023).

\bibitem{Haddara2023b}
M.~Haddara and E.~G.~Cavalcanti, New~J.~Phys. \textbf{25}, 093028 (2023).

\bibitem{Kent2023}
A.~Kent, arXiv:2302.12707 (2023).

\bibitem{Colladay1997}
D.~Colladay and V.~A.~Kosteleck\'{y}, Phys.~Rev.~D \textbf{55}, 6760 (1997).

\bibitem{Marsili2013}
F.~Marsili \emph{et~al.}, Nat.~Photonics \textbf{7}, 210--214 (2013).

\end{thebibliography}

\end{document}
